% !TEX root = ../main.tex
\begin{center}
    \vspace*{1.2cm}
    {\Large \textbf{Abstract}}
    \vspace{0.8cm}
\end{center}

\noindent
The study of viscoelasticity is a relevant problem across many fields and industries, 
both for its applications and its mathematical challenges.
Accurately modeling non-Newtonian behavior requires non-trivial governing
equations coupling momentum conservation with frame-indifferent constitutive
laws. To investigate this rheological complexity, the four-roll
mill provides a benchmark geometry: by modulating the relative rotation speeds
of four cylinders ($\Omega$ and $-k\Omega$, with $k \in [-1, 1]$), the apparatus
generates pure rotation, simple shear, pure planar elongation, and continuous
mixed kinematic regimes in between, providing an ideal platform to probe
viscoelastic response in these less ideal regimes.

\vspace{0.7em}
\noindent
However, characterizing viscoelastic properties from experimental velocity data
(such as Particle Image Velocimetry) presents an inverse problem that standard
Computational Fluid Dynamics (CFD) is not well suited to solve. Recovering unknown
constitutive parameters via mesh-based solvers requires an expensive, iterative
trial-and-error procedure involving repeated simulations. While Artificial Intelligence is
transforming engineering sciences, purely data-driven neural networks suffer
from heavy data requirements, lack of physical grounding, and poor extrapolation
outside training ranges. This thesis addresses these limitations by developing a
Physics-Informed Neural Network (PINN) framework applied to the four-roll mill.
By embedding Navier-Stokes momentum balance and viscoelastic constitutive laws
directly into the loss objective alongside data matching, the
optimization simultaneously obtain physical grounding, reduces training data
demands, and enables reliable field extrapolation.

\vspace{0.7em}
\noindent
Building upon foundational literature from Raissi's Newtonian flow past a cylinder
and Thakur's \textit{ViscoelasticNet} in open channels, this work expands
inverse viscoelastic characterization to an enclosed recirculating cavity with
zero net flux across an extended physical parameter range. The proposed
three-head architecture ($\mathcal{N}_\psi, \mathcal{N}_p, \mathcal{N}_{\boldsymbol{\tau}}$)
satisfies exact mass conservation via the stream function, algebraic pressure
anchoring, and a staged dual-precision optimization protocol (FP32 Adam + FP64
L-BFGS). 
This solver is capable of reconstructing unobserved stress and pressure fields
with $L_2$ errors below $1.0\%$ for velocity, $1.8\%$ for extra-stresses,
and $8.1\%$ for pressure. In the inverse problem, the network can 
simultaneously discover constitutive parameters and resolve spatial stress
fields for Oldroyd-B ($\lambda, \mu_p$ within $1.3\%$) and shear-thinning
Giesekus fluids ($\alpha, \lambda, \mu_p$ within $6.7\%$), maintaining robustness
down to 5,086 collocation points.
Finally, we map the intrinsic limits of recirculating kinematics,
revealing creeping-flow scaling invariance for solvent viscosity ($Re \ll 1$) 
and non-identifiability in linear PTT fluids, alongside proposals to overcome these limitations. 
\clearpage
